\section{Emergence, iterated}
\label{sec:iterated}

The bands of Figure~\ref{fig:rulers} are a family. They have members, the
bands; they have a count, which is the difference $D$ itself, rising by one
per band; and they have a profile, the tent of Figure~\ref{fig:tent}. So
they can beat. Lay a third ruler beside the first two, at a pitch such that
its count runs close to $D$, and a pattern of coincidences forms between
the third family and the \emph{bands}: a beat of a beat. This section says
when that happens, who can see it, and in what precise sense the pattern
that emerges is a new thing.

\begin{proposition}[Averages compose]
\label{prop:tower}
For two slides $w$ and $w'$, averaging along $w$ and then along $w'$ is the
same as averaging over both at once: it keeps exactly the recipes with
$k\cdot w=0$ and $k\cdot w'=0$. Pooling observers stacked in a cascade keep
smaller and smaller sets of recipes, and the set at each stage is the
intersection of the sets before it.
\end{proposition}

\begin{proof}
By Theorem~\ref{thm:sweep} the first average multiplies the wave of
recipe $k$ by $1$ if $k\cdot w=0$ and $0$ otherwise, and the second does the
same with $w'$; the product of the two indicators is the indicator of both
conditions at once.
\end{proof}

So a hierarchy of poolings is a hierarchy of recipe sets, and what a late
stage sees is a pattern in the counts the earlier stages left alone. For
three families the diagonal slide $(1,1,1)$ keeps every recipe whose
entries sum to zero: the three pairwise differences $(1,-1,0)$, $(0,1,-1)$,
$(1,0,-1)$ and, among others, $(1,-2,1)$, the difference of two differences.
That last recipe is the beat of beats, and whether an observer can see it
depends on what its front end can reach.

\begin{theorem}[Beats of beats need a front end of order four]
\label{thm:cascade}
Let three \emph{pure} tones --- sinusoids, each carrying only its first
harmonic --- add, with pairwise slow differences at rates $\delta_1$ and
$\delta_2$ nearly equal. The recipe $(1,-2,1)$, at rate $|\delta_1-\delta_2|$,
has order $|1|+|{-2}|+|1|=4$. A front end that is a polynomial of degree
$n$ mints from pure tones only recipes of order at most $n$, so a
square-law front end ($n=2$) and a cubic one ($n=3$) leave the beat of
beats silent. A cascade that squares, pools, and squares again reaches it,
because its first stage mints the two pairwise beats and its second
multiplies them; so does a single stage of degree four, and a hard
saturation, which is of every degree at once, reaches it weakly. A
multiplicative superposition of the three carries the recipe at linear
order, and a plain blur sees it.
\end{theorem}

\begin{proof}
A pure tone's wave has recipe $\pm1$ in its own count and $0$ in the
others (a real cosine is $\tfrac12(e^{2\pi i\theta}+e^{-2\pi i\theta})$, so
both signs are present). The $n$-th power of a sum of three such waves
expands into monomials $g_1^{e_1}g_2^{e_2}g_3^{e_3}$ with
$e_1+e_2+e_3=n$, and the recipes of such a monomial have $|k_i|\le e_i$,
hence order at most $n$; the recipe $(1,-2,1)$ needs $e_2\ge2$ and
$e_1,e_3\ge1$, four factors, which a square or a cube does not have. The
cascade's first square mints the pairwise differences $(1,-1,0)$ and
$(0,1,-1)$ at order two; the pooling keeps them and removes the rest; the
second square multiplies two real waves, and a product of two real waves
carries both the sum and the difference of their recipes --- the
difference here is $(1,-1,0)-(0,1,-1)=(1,-2,1)$, the beat of beats. A
product of three pictures carries every recipe by
Theorem~\ref{thm:whosees}.
\end{proof}

Measured, on three sinusoids at $300$, $330$ and $363$ hertz --- pairwise
beats at $30$ and $33$, beat of beats at $3$ --- as a fraction of a
first-order beat's strength: a square-law ear hears the three-hertz line
at $\earTernarySquare$, a cubic ear at $\earTernaryCubic$, the cascade at
$\earTernaryCascadePct\%$, and a multiplicative trio through a plain
low-pass at $\earTernaryPrintedPct\%$ (\texttt{ear.mjs};
Figure~\ref{fig:ternary}). The prediction for hearing is that second-order
beats of pure tones are heard only by a system whose front end reaches
order four --- two stages of detection (a front end followed by pooling),
or one stage of high order --- and the psychoacoustics of ``second-order
beats'' has data the theorem constrains; \S\ref{sec:ledger} names it as
the test.

\begin{figure}[t]
  \centering
  \begin{tikzpicture}
    \begin{axis}[width=0.7\linewidth, height=4.6cm, ybar, ymode=log,
      ymin=1e-8, ymax=3, bar width=16pt,
      ylabel={3 Hz line, as a fraction of a first-order beat},
      xtick={0,1,2,3}, xticklabels={square-law, cubic, {square, pool, square}, {multiplied, linear}},
      x tick label style={font=\scriptsize, align=center, text width=2.3cm},
      tick align=outside, tick pos=left, axis line style={gray!60},
      ymajorgrids, grid style={gray!18, very thin},
      label style={font=\footnotesize}, tick label style={font=\scriptsize},
      nodes near coords, point meta=rawy,
      every node near coord/.append style={font=\tiny, /pgf/number format/sci, /pgf/number format/sci generic={mantissa sep=\times,exponent={10^{#1}}}}]
      \addplot[fill=accent!70, draw=accent] table[x=index, y=fraction, col sep=comma]{data/ear-ternary.csv};
    \end{axis}
  \end{tikzpicture}
  \caption{\figlead{The beat of beats, and who hears it.} Three pure tones
  added in the air, listened to through a square-law ear, a cubic ear, and
  a two-stage ear that squares, pools and squares again; and the same three
  families multiplied, through a plain low-pass. The recipe $(1,-2,1)$ has
  order four: one stage of degree two or three cannot mint it from pure
  tones, two stages of degree two can, and a multiplicative superposition
  carries it from the start. Script \texttt{ear.mjs}, data
  \texttt{ear-ternary.csv}.}
  \label{fig:ternary}
\end{figure}

\begin{figure}[t]
  \centering
  \includegraphics[width=0.62\linewidth]{teaser-fantrio.png}
  \caption{\figlead{Three families, and what pooling sees of them.} Three
  fans of rays from three centres, drawn on the left half and pooled on the
  right. The diagonal slide keeps every recipe whose entries sum to zero,
  so the pooled half is a pattern in all three pairwise differences at
  once, and in their combinations: a rosette that no pair of the fans
  makes alone. Rendered by \Moire\ through its own pooled view
  (\S\ref{sec:instrument}).}
  \label{fig:trio}
\end{figure}

\subsection{What ``emergence'' means here}
\label{sec:emergence}

The word is used loosely often enough that it is worth saying exactly what
is claimed. The third pattern is a new family: it has members and a count,
and the count is a combined count of the parts. It is absent from every
linear description of the parts: an additive superposition has no cross
recipe, so no linear functional of the parts, no spectrum, no filter,
contains it. It appears only to observers that pool after a nonlinearity,
or in superpositions that multiply. And yet its geometry --- where its
members are, how fast its count runs, whether it exists at all --- is fixed
by the parts alone, through $\Phi$, and is the same for every observer that
can see it. That is emergence with a theorem attached, and the theorem is
the last one of this section. It is more abstract than the rest of the
paper; the reader who takes it on trust loses nothing that follows.

\begin{proposition}[The third pattern is the universal invariant]
\label{prop:invariant}
Fix a set of fast directions on the torus, and let $H$ be the set of all
slides along them, by any amount, whole or fractional: the
\emph{rephasings} of the carriers. Let $\E_HI$ be the picture averaged over
all of them. Call a rule $F$ that assigns a number to a picture
\emph{linear} if it respects sums and scalings (an integral of the picture
against any fixed weight is an example). Then every linear rule that gives
the same number for a picture and for every rephased copy of it,
$F(I(\cdot+h))=F(I)$ for all $h$ in $H$, is a function of the averaged
picture alone: $F(I)=F(\E_HI)$. And every rule applied to $\E_HI$ is
unchanged by rephasing. What can be measured of a superposition without
knowing the fast phases is exactly what can be measured of its third
pattern.
\end{proposition}

\begin{proof}
Average the invariance over $H$: $F(I)$ equals the average over $h$ of
$F(I(\cdot+h))$, and a linear rule passes through an average, so that
equals $F$ of the average of $I(\cdot+h)$, which is $F(\E_HI)$. The
converse is that $\E_HI$ is itself unchanged by any further rephasing along
$H$.
\end{proof}

The proposition says that ``the observable content of a superposition, to
anyone who cannot resolve the fast phases, is its tower of averages'' is
not a metaphor. It is the quotient of the picture by the fast rephasing,
in the sense in which a quotient is the most general thing invariant under
a symmetry: whatever two pictures differing only by a rephasing have in
common is a function of their common average, and nothing more. The word
\emph{linear} cannot be dropped --- the picture's energy $\int I^2$ is
unchanged by rephasing and is not a function of the average --- but the
observers that matter physically are not linear rules on the picture; they
are front ends followed by pooling, and for those
Theorem~\ref{thm:universality} already said that the picture is theirs and
the geometry is not. In the philosophers' terms this is \emph{weak}
emergence~\cite{Bedau1997Weak, Chalmers2006Emergence} --- the new pattern
is in principle derivable from the parts --- made exact: a new family, with
a new count, that no member of the parts' description contains, produced
by a quotient that is exact rather than approximate, and exactly computable
at every stage of the hierarchy. It is a better claim than strong
emergence made vague.

\begin{tryit}{a beat of beats}
Three oscillators on a synthesizer, sine waves, at $300$, $330$ and $363$
hertz. You will hear two fast throbs, at thirty and thirty-three per
second, and the question is whether you also hear a slow one, three times a
second. Theorem~\ref{thm:cascade} says a single mild nonlinearity cannot
produce it. Run the mix through a distortion effect and then a low-pass,
one stage: no three-hertz throb should appear. Run it through distortion,
a low-pass at about fifty hertz, and a second distortion: it should. If
your unaided ear hears the slow throb, that is data about the order of your
ear's front end.
\end{tryit}
